\documentclass[11pt,a4paper]{article}
\usepackage[utf8]{inputenc}
\usepackage[T1]{fontenc}
\usepackage[spanish,es-noshorthands]{babel}
\usepackage[margin=2.2cm]{geometry}
\usepackage{booktabs}
\usepackage{longtable}
\usepackage{array}
\usepackage{xcolor}
\usepackage{enumitem}
\usepackage{hyperref}
\usepackage{fancyhdr}

\definecolor{azul}{RGB}{28,56,121}
\hypersetup{colorlinks=true,linkcolor=azul,urlcolor=azul}
\pagestyle{fancy}
\fancyhf{}
\fancyhead[L]{\small TEL354 \textperiodcentered{} Grupo 3 (H0892) \textperiodcentered{} AVZ9 \textperiodcentered{} Consolidación del LLSD de R1, R2 y RE}
\fancyhead[R]{\small \thepage}

\begin{document}

\begin{titlepage}
\centering
\vspace*{1.2cm}
{\Large\bfseries Entregable Semana 9 (AVZ9)}\\[0.6cm]
{\LARGE\bfseries\color{azul} Consolidación del LLSD de R1, R2 y RE\\[0.2cm] con análisis cuantitativo}\\[0.4cm]
{\large Segunda parte}\\[1.2cm]
{\large Grupo 3 --- Horario H0892}\\[0.2cm]
{\large TEL354: Redes Definidas por Software --- 2026-2}\\[0.2cm]
{\large 31 de octubre de 2026}\\[1.6cm]
\begin{tabbing}
Coach \hspace{1.4cm} \= Germán Espinoza\\[2pt]
Profesor \> Christian Isaac Quispe Ordoñez\\[2pt]
Integrantes \> Rojas Soplin, Rodrigo Franco --- 20213686 (líder)\\[2pt]
\> Oropeza León, Piero Jesús --- 20212733 (arquitecto)\\[2pt]
\> Altamirano Ramírez, Camila Nicole --- 20220964 (PM / QA)\\[2pt]
\> Tavella Rivero, Fernando Vincenzo --- 20182693 (DevOps / SRE)\\[6pt]
Archivo Paideia \> AVZ9-TEL354\_2026-2\_G3
\end{tabbing}
\end{titlepage}

\section*{Cómo leer este documento}
Este documento consolida el diseño de bajo nivel de la solución para los requisitos \textbf{R1} (control de acceso según rol), \textbf{R2} (protección de recursos privilegiados) y \textbf{RE} (eficiencia), y lo acompaña con su análisis cuantitativo: cada afirmación de diseño tiene su cifra, y cada cifra dice de dónde salió. Está escrito para leerse sin conocer los avances anteriores: la sección 2 presenta la solución completa, las secciones 3 a 21 detallan el diseño pieza por pieza ---algoritmos, estados, contratos, matriz de acceso, configuración y secuencias completas---, las secciones 22 a 24 presentan los números, y la sección 25 los contrasta con el avance previo justificando cada ajuste.

\section{La solución en una página}

La solución es una red de campus gobernada por un controlador SDN, donde \emph{quién puede qué} se decide en un punto central y se materializa como reglas concretas en los switches. El flujo de vida de una decisión es siempre el mismo: un dispositivo se conecta y recibe un \textbf{perfil mínimo (BASE)}; para salir de ese mínimo, la persona se autentica en un \textbf{portal cautivo}; el sistema de identidad determina \emph{quién es} y \emph{con qué rol}; el motor de políticas decide \emph{qué puede}; el adaptador del controlador traduce la decisión a \textbf{reglas OpenFlow} con prioridades fijas; el switch las ejecuta y cuenta; y la auditoría registra todo.

Tres realidades definen el diseño. Primera: existen \textbf{tres reinos de identidad} --- la comunidad universitaria se autentica contra el proveedor de identidad institucional; los operadores (TI, administradores, superadministradores) contra un repositorio propio de la plataforma con segundo factor obligatorio; y los invitados contra un registro efímero que se borra solo al vencer. Segunda: el tráfico se gobierna por una \textbf{escalera de prioridades} que es una jerarquía de autoridad --- la mitigación manda sobre la sesión, la sesión sobre el perfil base, y el perfil base sobre los caminos. Tercera: la red se instala \textbf{de forma proactiva} --- lo que se sabe que va a ocurrir (el perfil base, la sesión, la mitigación) se escribe en el switch de inmediato; el controlador solo se consulta para las tres excepciones cerradas: aprender un dispositivo nuevo, responder un ARP sin resolver y descubrir la topología.

\begin{longtable}{p{2.8cm}p{7.2cm}p{4.4cm}}
\toprule
\textbf{Componente} & \textbf{Qué hace} & \textbf{Se comunica por}\\
\midrule
Portal cautivo & La única interfaz que ven las personas: elige el reino y pide credenciales; no decide nada & HTTP (personas); API REST (identidad)\\
Servicio de identidad & Verifica contra el reino elegido, emite la sesión con su rol y su vigencia, y la retira & REST \textperiodcentered{} RADIUS \textperiodcentered{} bus de eventos\\
Registro de dispositivos & Guarda qué equipos están habilitados y su historial; registrar no concede privilegios & REST \textperiodcentered{} bus de eventos\\
Motor de políticas & El único punto de decisión: evalúa identidad, contexto y recurso, y ordena & REST \textperiodcentered{} bus de eventos\\
Adaptador del controlador & Traduce cada orden a reglas OpenFlow con su prioridad, cookie y vencimiento & REST (northbound)\\
Controlador SDN & Clúster de tres nodos; mantiene la topología y aplica las reglas & OpenFlow (southbound)\\
Switches & Ejecutan y cuentan: conmutan sin depender del controlador & OpenFlow\\
Bus de eventos & Lo asíncrono: informa, no ejecuta & AMQP con colas durables\\
Auditoría & Registra todo lo que ocurre, sin excepción & Bus (escucha el catálogo completo)\\
\bottomrule
\end{longtable}

\section{Evolución del diseño: qué consolidó esta semana}

El avance anterior dejó planteadas once cuestiones abiertas. Esta sección muestra cómo quedaron resueltas; es la evolución del diseño, no su repetición, y cada resolución cita la evidencia que la sostiene.

\begin{longtable}{p{4.2cm}p{5.2cm}p{4.6cm}}
\toprule
\textbf{Cuestión abierta del avance previo} & \textbf{Decisión consolidada} & \textbf{Evidencia}\\
\midrule
¿El switch soporta múltiples tablas y metadatos entre ellas? & Sí: 254 tablas disponibles (0--253). El pipeline de dos etapas puede materializarse en tablas físicas separadas & Sonda directa sobre el switch del laboratorio\\
¿Cuál es la capacidad real de TCAM del switch? & No hay TCAM física: la conmutación es por software. El límite real de entradas se mide por llenado progresivo & Sonda; presupuesto en §24\\
¿Cómo se redirige al portal? & La entrada de prioridad 150 del perfil base lleva al portal; el dispositivo sin sesión solo alcanza el portal, DHCP y DNS & Diseño del pipeline (§17)\\
¿Qué tasa de reglas por segundo y qué latencia de confirmación soporta el controlador? & Toda orden se confirma por barrera OpenFlow antes de responder; la medición fina es parte de la campaña de pruebas & Implementación del adaptador (§11)\\
¿Quién es elegible para perfiles elevados? & Solo la autenticación o una elevación aprobada conceden privilegios; no hay elegibilidad implícita & Política de autorización (§7)\\
¿El directorio viaja en claro dentro del segmento? & En el prototipo se acepta y se documenta; en despliegue, cifrado obligatorio & Decisión documentada\\
¿La suplantación de MAC e IP en el mismo puerto es detectable? & No por el pipeline; se declara como limitación conocida y se compensa con la escalera & Análisis de amenazas\\
¿Cómo agregan caminos los switches del núcleo? & Prefijos por switch de acceso en la referencia; caminos individuales en el prototipo & Diseño de enrutamiento\\
¿Clúster del controlador de dos o de tres nodos? & Tres nodos: el consenso exige mayoría, y tres es el mínimo que sobrevive a una caída & Decisión de alta disponibilidad\\
¿Cuántos recursos requiere el despliegue? & 14 máquinas virtuales: 21 núcleos, 31 GB de RAM y 145 GB de disco, repartidos por rol & Tabla §15\\
¿Quedan inconsistencias con la propuesta presentada? & La propuesta consolidada ya refleja las prioridades y las cookies definitivas & Documento de propuesta\\
\bottomrule
\end{longtable}

\section{Arquitectura interna de los servicios de R1 y R2}

Hay un solo punto de decisión ---el motor de políticas--- y un solo traductor a OpenFlow ---el adaptador del controlador---. El resto identifica, aporta atributos o presenta. Los servicios no se llaman entre sí por invocaciones internas: lo síncrono viaja por REST y lo que informa, por el bus de eventos.

\subsection{Portal cautivo}
Sin estado propio. Presenta el selector de reino ---comunidad, operación o invitado--- y reenvía las credenciales al servicio de identidad. El selector \emph{no concede nada}: solo elige la ruta de verificación; un intento de operador con credenciales de comunidad ---o al revés--- falla en la verificación y queda auditado. La página cautiva es una interfaz sobre su propia API: entrada, registro de invitado, solicitud de elevación y cierre de sesión.

\subsection{Servicio de identidad}
Dueño de las sesiones y de los repositorios de credenciales de operadores e invitados. Su algoritmo de verificación, reino por reino, está en §5; su máquina de estados de sesión, en §6.

\subsection{Registro de dispositivos}
Guarda los equipos habilitados y su historial (altas, revocaciones, movimientos). Su regla de fondo: \emph{registrar no concede confianza} --- el registro es un atributo que la política evalúa, jamás un privilegio en sí mismo. Solo los operadores se registran; los usuarios académicos no (mantiene el registro acotado).

\subsection{Motor de políticas}
El único componente que decide. Evalúa la asignación \emph{rol + permiso + recurso + condiciones} contra la política central (las once condiciones de §7) y emite la orden correspondiente: aplicar perfil de sesión, aplicar mitigación o retirar. Toda decisión de mitigación lleva vencimiento obligatorio ---ninguna regla de respuesta queda instalada para siempre--- y las decisiones de alto impacto exigen aprobación humana previa.

\subsection{Adaptador del controlador}
Traduce cada orden a reglas OpenFlow concretas. Es el único componente que conoce la sintaxis del controlador; los demás expresan \emph{qué}, nunca \emph{cómo}. Confirma cada instalación con una barrera OpenFlow ---la regla está escrita en el switch antes de responder--- y retira \emph{siempre por cookie}: nunca toca una regla ajena.

\subsection{Consola y Auditoría}
La consola es la interfaz del operador: alertas en vivo y aprobaciones. Su sesión es la del propio operador (efecto SSO: una sola autenticación abre lo que el rol permite) y cada petición se valida contra el rol en el servidor. La auditoría escucha el catálogo completo de eventos y registra también los huecos: si un evento llega con número de secuencia saltado, el hueco queda \emph{declarado}, jamás oculto.

\section{El algoritmo del login, reino por reino}

\subsection{Comunidad universitaria (proveedor institucional por RADIUS)}
\begin{enumerate}
\item La persona elige «comunidad» en el selector y envía usuario y contraseña.
\item El portal reenvía las credenciales al servicio de identidad, sin guardarlas.
\item El servicio de identidad, como cliente RADIUS, envía una \emph{petición de acceso} al servidor RADIUS institucional (usuario y contraseña cifrados según el protocolo).
\item El servidor responde \emph{acceso aceptado} o \emph{acceso rechazado}. Con rechazo, el intento queda auditado y termina el flujo.
\item Con aceptación, el servicio abre el registro contable de la sesión (\emph{inicio}) y crea la sesión con perfil académico.
\item La sesión se publica como evento y se ordena al controlador la instalación del perfil (prioridad 200, §17).
\end{enumerate}
El prototipo simula al proveedor institucional con un servidor RADIUS propio y un directorio de cuentas de prueba: se simula \emph{al proveedor}, no \emph{al protocolo}.

\subsection{Operadores (repositorio propio + TOTP)}
\begin{enumerate}
\item El aprovisionamiento ocurre antes del primer login: un administrador crea la cuenta (usuario, contraseña con hash bcrypt, rol) y el sistema genera el secreto del segundo factor, mostrado una sola vez como código QR.
\item El operador carga el QR en su aplicación autenticadora y verifica un código; recién entonces el par (dispositivo, operador) queda habilitado.
\item En el login: usuario y contraseña se verifican contra el hash bcrypt ---jamás contra un texto en claro---.
\item El código TOTP del momento se verifica contra el secreto. Sin código correcto no hay sesión, y no existe variante sin segundo factor.
\item La sesión se emite con el rol del operador (TI, administrador o superadministrador) y sus relojes más estrictos (§6).
\end{enumerate}

\subsection{Invitados (registro efímero)}
\begin{enumerate}
\item El invitado se autoregistra con usuario y contraseña; la cuenta nace con vencimiento a 24 horas.
\item El login verifica contra ese registro como cualquier otro reino.
\item Al vencer el plazo, la cuenta se borra sola --- sin intervención administrativa --- y los intentos posteriores se rechazan.
\end{enumerate}
El mecanismo completo del reino invitado está fuera del alcance del prototipo; el lugar queda reservado y operativo en lo esencial.

\section{La sesión: estados y relojes}

\begin{longtable}{p{3.2cm}p{5.4cm}p{5.4cm}}
\toprule
\textbf{Estado} & \textbf{Cómo se entra} & \textbf{Cómo se sale}\\
\midrule
Activa & Login verificado: sesión emitida con rol, dispositivo y relojes & Cierre explícito \textperiodcentered{} revocación administrativa \textperiodcentered{} inactividad \textperiodcentered{} tope \textperiodcentered{} movimiento del dispositivo\\
Cerrada & Cualquiera de las salidas anteriores & Terminal; el evento \emph{sesión cerrada} registra el motivo\\
\bottomrule
\end{longtable}

Los dos relojes, por perfil:

\begin{longtable}{p{4.4cm}p{3.0cm}p{3.0cm}p{3.6cm}}
\toprule
\textbf{Perfil} & \textbf{Inactividad máxima} & \textbf{Tope absoluto} & \textbf{Sentido}\\
\midrule
Académico & 30 min & 12 h & Acompaña la jornada; la inactividad devuelve el equipo al perfil base\\
Operador (TI, Admin, SA) & 15 min & 8 h & Más estricto: la sesión abre la red de gestión y la consola\\
Elevación temporal & según la elevación & 1 h por defecto & El tope es parte de la decisión que la concede\\
\bottomrule
\end{longtable}

La misma cifra vive en dos lugares ---la regla del switch y la sesión--- y es una sola por perfil: sin desincronización posible. La sesión está además \emph{ligada al dispositivo}: nace atada a la ubicación (MAC, puerto, switch, IP) que el controlador aprendió, y si esa MAC aparece en otro puerto o switch, la sesión se invalida de inmediato. Un cierre explícito no espera a ningún reloj.

\section{El motor de políticas y sus condiciones}

La política central define 4 roles, 37 permisos y 78 asignaciones. Cada asignación dice: \emph{este rol, sobre este permiso, para estos recursos, si se cumplen estas condiciones, con estas obligaciones} (auditar, notificar, justificar, registrar). Las once condiciones son la letra fina del «quién puede qué»:

\begin{longtable}{p{1.6cm}p{4.6cm}p{7.8cm}}
\toprule
\textbf{Cond.} & \textbf{Nombre} & \textbf{Qué exige}\\
\midrule
C1 & Ámbito & El recurso pertenece al ámbito asignado al rol\\
C2 & Delegación vigente & Existe una delegación explícita, con alcance y vencimiento\\
C3 & Rango aprobado & El valor o la acción está dentro del rango aprobado\\
C4 & Separación de funciones & El solicitante es distinto del aprobador\\
C5 & Sesión privilegiada & Segundo factor verificado, dispositivo registrado y canal de gestión\\
C6 & Estado de seguridad & Sin compromiso, cuarentena ni incidente que impida la acción\\
C7 & Bajo impacto & La severidad no supera el umbral; si lo supera, se exige aprobación\\
C8 & Sesión propia & El sujeto solo opera sus propias sesiones\\
C9 & Recurso autorizado & El tipo y la instancia del recurso están en la lista del rol\\
C10 & Contexto permitido & Ventana de mantenimiento o emergencia declarada y ubicación permitida\\
C11 & Solo lectura & El rol consulta y ejecuta operaciones autorizadas, no configura\\
\bottomrule
\end{longtable}

Reglas transversales que nadie puede violar: nadie aprueba su propia solicitud; nadie se asigna un rol a sí mismo; la combinación de roles no supera el mayor privilegio (la delegación es la única vía de ampliación); solo el superadministrador delega; y la auditoría nunca la administra quien opera. La política es un documento único, validado automáticamente: toda referencia cruzada (rol, permiso, condición, recurso) se comprueba en cada cambio, y su vista de matriz se regenera de la misma fuente.

\section{Modelo de datos}

\begin{longtable}{p{3.4cm}p{6.0cm}p{4.6cm}}
\toprule
\textbf{Estructura} & \textbf{Campos} & \textbf{Notas}\\
\midrule
Sesión & identificador opaco \textperiodcentered{} identidad \textperiodcentered{} rol o perfil \textperiodcentered{} MAC \textperiodcentered{} puerto \textperiodcentered{} switch \textperiodcentered{} IP \textperiodcentered{} inactividad máxima \textperiodcentered{} tope absoluto & Vive en el servicio de identidad; el retiro es inmediato\\
Operador & usuario \textperiodcentered{} hash bcrypt \textperiodcentered{} secreto TOTP \textperiodcentered{} rol & Creado por la cadena SA$
ightarrow$Admin$
ightarrow$TI; el secreto se muestra una vez\\
Invitado & usuario \textperiodcentered{} hash bcrypt \textperiodcentered{} vencimiento & Se borra solo al vencer (24 h)\\
Dispositivo registrado & MAC \textperiodcentered{} titular \textperiodcentered{} rol \textperiodcentered{} responsable \textperiodcentered{} vigencia \textperiodcentered{} historial & Registrar no concede confianza\\
Decisión de política & incidente \textperiodcentered{} acción \textperiodcentered{} parámetros (tasa, TTL) \textperiodcentered{} resultado & Toda mitigación expira sola\\
Evento & identificador único \textperiodcentered{} tipo \textperiodcentered{} emisor \textperiodcentered{} marca de tiempo \textperiodcentered{} secuencia \textperiodcentered{} correlación \textperiodcentered{} contenido & La envoltura es común a los dieciséis tipos del catálogo\\
\bottomrule
\end{longtable}

\section{Contratos de interfaz, campo a campo}

\subsection{Órdenes al controlador (northbound)}
Tres órdenes, un mismo recurso, confirmación en dos niveles (aplicada + verificación diferida):

\begin{longtable}{p{3.0cm}p{3.0cm}p{8.0cm}}
\toprule
\textbf{Orden} & \textbf{Campos} & \textbf{Semántica}\\
\midrule
Aplicar perfil & dispositivo (MAC, switch, puerto) \textperiodcentered{} vigencia (inactividad y tope) \textperiodcentered{} contexto (cookie, solicitante, motivo) \textperiodcentered{} perfil & Instala la sesión con prioridad 200; la vigencia es la misma cifra en la regla y en la sesión\\
Aplicar mitigación & acción (limitar, bloquear origen, aislar, cuarentena) \textperiodcentered{} objetivo \textperiodcentered{} orígenes \textperiodcentered{} parámetros (tasa, TTL) \textperiodcentered{} contexto (cookie, incidente, solicitante, aprobación) & Instala la respuesta con prioridad 1000 y vencimiento = TTL; la cuarentena se rechaza sin aprobación\\
Retirar & contexto (cookie) & Retiro por cookie, siempre; responde al confirmar\\
\bottomrule
\end{longtable}

Reglas de las órdenes: la cookie es la clave de idempotencia (reaplicar la misma cookie no duplica reglas); ante un controlador que no responde, el solicitante reintenta \emph{la misma orden}; los errores son estructurados (400 esquema inválido, 401 sin clave, 403 fuera de competencia, 404 desconocido, 409 duplicado, 503 no disponible).

\subsection{Eventos del bus}
Todo evento viaja con la misma envoltura: identificador único (clave de idempotencia), tipo, emisor, marca de tiempo, secuencia por incidente y claves de correlación. El contenido es propio de cada tipo:

\begin{longtable}{p{4.2cm}p{5.2cm}p{4.6cm}}
\toprule
\textbf{Evento} & \textbf{Campos del contenido} & \textbf{Quién produce $
ightarrow$ quién consume}\\
\midrule
Sesión abierta & identidad \textperiodcentered{} rol o perfil \textperiodcentered{} dispositivo \textperiodcentered{} marca de tiempo & Identidad $
ightarrow$ políticas, auditoría\\
Sesión cerrada & identidad \textperiodcentered{} rol o perfil \textperiodcentered{} dispositivo \textperiodcentered{} motivo & Identidad $
ightarrow$ políticas, auditoría\\
Dispositivo registrado / revocado & MAC \textperiodcentered{} titular \textperiodcentered{} rol \textperiodcentered{} responsable \textperiodcentered{} vigencia & Registro $
ightarrow$ auditoría, consola\\
Solicitud de elevación & solicitante \textperiodcentered{} perfil pedido \textperiodcentered{} alcance \textperiodcentered{} duración & Portal $
ightarrow$ políticas, consola, auditoría\\
Mitigación requerida & incidente \textperiodcentered{} acción \textperiodcentered{} objetivo \textperiodcentered{} orígenes \textperiodcentered{} parámetros \textperiodcentered{} TTL & Políticas $
ightarrow$ comando al controlador; eco al bus\\
Mitigación aplicada & incidente \textperiodcentered{} reglas \textperiodcentered{} switches \textperiodcentered{} marca de tiempo & Controlador $
ightarrow$ monitor, incidentes, auditoría\\
Mitigación verificada & incidente \textperiodcentered{} tasas antes/después \textperiodcentered{} cumple objetivo & Monitor $
ightarrow$ incidentes, políticas, auditoría\\
Mitigación expirada & incidente \textperiodcentered{} regla \textperiodcentered{} motivo & Monitor/controlador $
ightarrow$ incidentes, políticas, auditoría\\
\bottomrule
\end{longtable}

La regla de fondo del bus: \textbf{el evento informa; el comando ejecuta}. Las decisiones que deben aplicarse con confirmación viajan como comandos por la API del controlador y publican su eco por el bus para que la auditoría y la consola las vean. Nunca viajan credenciales, ni en el contenido ni en la correlación. La entrega es al menos una vez, con idempotencia por identificador y orden por incidente; lo que no se procesa tras cinco reintentos queda en una cola de no entregados, visible y auditada.

\section{El pipeline del switch}

La regla es la unidad de política materializada: cada entrada del switch es una decisión convertida en coincidencia y acción. El orden de evaluación es una jerarquía de autoridad ---a mayor prioridad, mayor poder de decidir sobre el paquete--- y es el mismo en toda la red:

\begin{longtable}{p{1.8cm}p{4.6cm}p{3.4cm}p{2.2cm}p{2.0cm}}
\toprule
\textbf{Prior.} & \textbf{Entrada} & \textbf{Acción} & \textbf{Cookie} & \textbf{Vigencia}\\
\midrule
1000 & Mitigación (limitar tasa) & medidor con banda de descarte, luego caminos & incidente & TTL de la mitigación\\
1000 & Mitigación (bloquear origen / aislar / cuarentena) & descarte & incidente & TTL corto o de la aprobación\\
200 & Sesión de perfil & hacia el alcance del perfil & sesión & inactividad y tope de la sesión\\
200 & Elevación & hacia el alcance elevado & sesión temporal & TTL de la elevación\\
150 & Portal (pre-login) & hacia el portal & esqueleto & permanente\\
120 & Par MAC/IP aprendido & hacia caminos & esqueleto & permanente\\
110 & Anti-suplantación (IP con MAC ajena) & descarte & esqueleto & permanente\\
100 & Destino a la red de gestión & descarte & esqueleto & permanente\\
20 & Caminos & puerto calculado & camino & mientras valga la topología\\
10 & DHCP y DNS & hacia caminos & esqueleto & permanente\\
--- & Sin coincidencia & paquete al controlador (aprendizaje) & --- & ---\\
\bottomrule
\end{longtable}

Tres lecturas de esta tabla, para que quede clara la jerarquía:

\begin{itemize}
\item \textbf{La mitigación manda sobre todo} --- es la respuesta al ataque y no puede ser vencida por una sesión. \textbf{La sesión manda sobre el esqueleto} --- es el privilegio otorgado. \textbf{El esqueleto manda sobre los caminos} --- es la línea base de cualquier dispositivo, autenticado o no.
\item \textbf{El esqueleto (prioridades 150 a 10) es la lista cerrada} que todo dispositivo lleva antes de autenticarse: portal, par MAC/IP, anti-suplantación, denegación a la gestión y DHCP/DNS. La sesión no lo reemplaza: convive con él y lo vence por prioridad. El descarte hacia la red de gestión es permanente --- una sesión académica jamás lo abre.
\item \textbf{El modo de fallo es seguro.} Si el controlador deja de responder, las entradas instaladas permanecen y un paquete sin coincidencia se descarta: el switch no vuelve a conmutación normal ni inunda. La red no se abre por la caída del controlador; se degrada a la última política conocida. El clúster de tres nodos cubre la continuidad del control; el modo seguro cubre el intervalo.
\end{itemize}

Las cookies codifican tipo y correlativo en 64 bits: los bits altos distinguen sesión, incidente, esqueleto y camino. El controlador mantiene el registro entre la clave legible y el valor, y el retiro es siempre por cookie, en todas las tablas: nada queda pegado.

\subsection{Ejemplo con números: una docente en el puerto p20}
Una docente conecta su equipo en el puerto p20 del switch de acceso A1. Sucede, en orden:

\begin{enumerate}
\item \textbf{Aprendizaje (una sola vez).} El primer paquete no coincide con nada y sube al controlador, que aprende el par (MAC, IP) y su ubicación (A1, p20). Desde ese momento, las entradas proactivas gobiernan: el equipo alcanza el portal (150), recibe DHCP y DNS (10), su tráfico con otra MAC es descartado (110) y todo destino hacia la red de gestión es descartado (100).
\item \textbf{Login.} La docente entra al portal y se autentica contra el proveedor institucional. El servicio de identidad emite su sesión con perfil académico (inactividad 30 min, tope 12 h).
\item \textbf{Instalación de la sesión.} Llega la orden al controlador: instalar en A1 la regla de prioridad 200 para (MAC, IP) con cookie de sesión y vigencia igual a la sesión. Desde ahora, el tráfico de la docente hacia los servicios de la comunidad e Internet pasa por la prioridad 200.
\item \textbf{La escalera decide.} Un paquete hacia la gestión sigue cayendo en el descarte de prioridad 100 --- la sesión no lo abre. Un paquete hacia un servicio académico coincide con la sesión (200) y sigue a los caminos (20). Si mañana hubiera una mitigación contra su equipo, la prioridad 1000 vencería a la sesión.
\item \textbf{Cierre.} Al cerrar sesión ---o al vencer un reloj---, el retiro es por cookie: solo las reglas de esa sesión desaparecen; el esqueleto permanece.
\end{enumerate}

\section{El adaptador del controlador por dentro}

El adaptador es el único componente que conoce la sintaxis del controlador. Su interior:

\begin{longtable}{p{4.6cm}p{9.4cm}}
\toprule
\textbf{Pieza} & \textbf{Qué hace, en detalle}\\
\midrule
Registro de cookies & Traduce cada clave legible (sesión, incidente) a su cookie de 64 bits: los bits altos distinguen el tipo (0x01 sesión, 0x02 incidente, 0x03 esqueleto, 0x04 camino) y el resto es correlativo. El retiro se hace siempre por esta vía: se recorren las reglas del switch y solo se retiran las que llevan la cookie propia\\
Aplicar mitigación & Por acción, construye la regla: \emph{limitar tasa} instala primero un medidor (unidad: paquetes por segundo, banda de descarte a la tasa de la decisión) y luego la regla de prioridad 1000 que casa origen, puerto y destino y aplica el medidor antes de los caminos; \emph{bloquear origen}, \emph{aislar} y \emph{cuarentena} instalan la regla de prioridad 1000 con acción de descarte. Todas llevan vencimiento igual al TTL de la decisión. Si el origen no está identificado, solo se permite la limitación por destino ---aplicada en todos los bordes---; bloquear sin origen se rechaza: bloquearía todo el borde\\
Aplicar perfil & Regla de prioridad 200 que casa el puerto y la MAC del dispositivo y entrega al camino; lleva los dos relojes de la sesión (inactividad y tope) con la misma cifra que la sesión\\
Confirmación & Cada instalación se confirma con una barrera OpenFlow: la orden responde \emph{aplicada} con las reglas creadas y la marca de tiempo, y la verificación (si aplica) queda pendiente para el monitor\\
Consultas & Asociación vigente (MAC, puerto, switch, IP) desde los anfitriones conocidos; reglas instaladas por cookie; topología y contadores del controlador; e inventario de servicios declarados\\
\bottomrule
\end{longtable}

\section{La matriz de acceso: perfiles y recursos}

El plano de red se gobierna por perfiles. La escalera de perfiles, con su prioridad y su alcance:

\begin{longtable}{p{3.6cm}p{1.8cm}p{3.0cm}p{5.6cm}}
\toprule
\textbf{Perfil} & \textbf{Prior.} & \textbf{Cuándo} & \textbf{Alcance}\\
\midrule
BASE & 150/10 & Siempre (pre-login) & Portal, DHCP y DNS --- la lista cerrada\\
ACADÉMICO & 190 & Tras login de comunidad & Servicios de la comunidad e Internet\\
OPERADOR\_TI & 200 & Tras login de operador & Consola y API de seguridad\\
OPERADOR\_ADMIN & 200 & Tras login de administrador & Red de gestión\\
OPERADOR\_SA & 200 & Tras login de superadministrador & Red de gestión y acceso remoto\\
ELEVACIÓN & 200 & Aprobada, con vencimiento & Según la instancia (p. ej., laboratorio)\\
\bottomrule
\end{longtable}

La matriz perfil $\times$ recurso resultante (plano de red):

\begin{longtable}{p{3.4cm}p{2.2cm}p{2.4cm}p{2.4cm}p{2.4cm}p{2.2cm}}
\toprule
\textbf{Recurso} & \textbf{BASE} & \textbf{ACAD.} & \textbf{OP. TI} & \textbf{OP. ADMIN} & \textbf{OP. SA}\\
\midrule
Portal & SÍ & SÍ & SÍ & SÍ & SÍ\\
DHCP y DNS & SÍ & SÍ & SÍ & SÍ & SÍ\\
Servicios de la comunidad & --- & SÍ & SÍ & SÍ & SÍ\\
Internet & --- & SÍ & SÍ & SÍ & SÍ\\
Consola & --- & --- & SÍ & SÍ & SÍ\\
API de seguridad & --- & --- & SÍ & SÍ & SÍ\\
API de red & --- & --- & --- & SÍ & SÍ\\
Red de gestión & --- & --- & --- & SÍ & SÍ\\
Acceso remoto & --- & --- & --- & --- & SÍ\\
Laboratorio (elevación) & --- & --- & --- & --- & ---\\
\bottomrule
\end{longtable}

La columna de laboratorio es intencionalmente vacía: el laboratorio no es un perfil fijo, es una instancia de elevación con su solicitante, su alcance y su vencimiento. En el plano de servicios, la matriz equivalente es la de las 78 asignaciones rol $\times$ permiso de la política (§7); aquí se muestra el plano de red, que es el que el switch materializa.

\section{Los endpoints del portal (I3), campo a campo}

\begin{longtable}{p{3.0cm}p{1.6cm}p{4.4cm}p{3.0cm}p{2.0cm}}
\toprule
\textbf{Recurso} & \textbf{Método} & \textbf{Campos de entrada} & \textbf{Respuesta} & \textbf{Errores}\\
\midrule
/login & GET & --- & Selector con las tres poblaciones & ---\\
/login & POST & población \textperiodcentered{} usuario \textperiodcentered{} contraseña \textperiodcentered{} TOTP (operación) & Sesión \{identificador, perfil\} & 401, 503\\
/logout & POST & identificador de sesión & \{estado: cerrada\} & 404, 503\\
/guest-register & POST & usuario \textperiodcentered{} contraseña & \{estado, vencimiento\} & 503\\
/elevation & GET & --- & Alcances disponibles & ---\\
/elevation & POST & solicitante \textperiodcentered{} alcance \textperiodcentered{} duración & 202 \{estado, identificador de evento\} & 503\\
\bottomrule
\end{longtable}

Reglas del portal: la contraseña nunca se guarda ni se registra; la población elegida es un dato de auditoría del intento; un intento con credenciales en la ruta equivocada falla y queda auditado.

\section{Los endpoints de administración (I7), campo a campo}

Cada servicio expone su administración; la consola la agrega y la sesión del operador es la autenticación única. Cada petición se valida contra la sesión \emph{y} el rol en el servidor:

\begin{longtable}{p{3.0cm}p{4.6cm}p{1.8cm}p{4.6cm}}
\toprule
\textbf{Servicio} & \textbf{Recurso} & \textbf{Método} & \textbf{Rol mínimo (permiso)}\\
\midrule
Identidad & /sessions/\{id\} & GET \textperiodcentered{} DELETE & TI (sesiones propias o administrativas)\\
Identidad & /operators & POST & Administrador (cuentas privilegiadas)\\
Registro & /devices \textperiodcentered{} /devices/\{mac\}/history & GET \textperiodcentered{} POST & TI / Administrador\\
Políticas & /profiles & GET \textperiodcentered{} PUT & PUT: Administrador (P17)\\
Políticas & /thresholds & GET \textperiodcentered{} PUT & PUT: TI / Admin / SA (P33)\\
Políticas & /elevations & GET & TI / Administrador\\
Auditoría & /events (desde, actor, cursor) & GET & Solo lectura\\
Monitor & /metrics \textperiodcentered{} /baselines & GET & Solo lectura\\
Consola & /alerts & GET & TI / Admin / SA (P09)\\
\bottomrule
\end{longtable}

Convenciones comunes: respuesta \{datos, error\{código, mensaje\}\}; paginación por cursor; 401 sin sesión válida, 403 con sesión sin el rol. Quien cambia perfiles o umbrales queda registrado con su identidad.

\section{Configuración de RADIUS y LDAP}

El servicio de identidad es cliente RADIUS del proveedor institucional. La configuración del prototipo:

\begin{longtable}{p{5.0cm}p{9.0cm}}
\toprule
\textbf{Parámetro} & \textbf{Valor y sentido}\\
\midrule
Servidor RADIUS & El servidor del prototipo (simula al proveedor institucional; en despliegue, el institucional)\\
Secreto compartido & Clave entre el cliente (identidad) y el servidor; jamás en claro en configuraciones versionadas\\
Red de clientes autorizados & 10.0.0.0/24: solo la identidad puede consultar\\
Contabilidad & Identificador contable único por sesión; \emph{inicio} al abrir, \emph{intermedio} periódico y \emph{fin} al cerrar --- alimenta la correlación de la auditoría\\
Directorio detrás de RADIUS & Cuentas de prueba del proveedor simulado; el prototipo simula al proveedor, no al protocolo\\
Transporte & Prototipo: 1812/1813 y directorio en claro, documentado; despliegue: cifrado obligatorio\\
\bottomrule
\end{longtable}

\section{Secuencia: el esqueleto al conectar (proactivo)}

Lo que ocurre cuando un switch se incorpora --- sin que ningún paquete de usuario haya aparecido todavía:

\begin{enumerate}
\item El switch se conecta al controlador y negocia OpenFlow 1.3; el controlador confirma sus capacidades (254 tablas, medidores, grupos).
\item El controlador instala de inmediato el esqueleto estático: portal (150), denegación a la gestión (100), DHCP/DNS (10) y los caminos conocidos (20). Nada de esto espera a un usuario.
\item Llega el primer paquete de un dispositivo nuevo: sin coincidencia, sube al controlador --- es la única vez.
\item El controlador aprende el par (MAC, IP) y su ubicación, y completa el esqueleto: par MAC/IP (120) y anti-suplantación (110).
\item Desde ese instante, todo el tráfico del dispositivo coincide con reglas proactivas; el controlador no vuelve a estar en el camino.
\end{enumerate}

\section{Secuencia: la mitigación completa}

El caso de ataque más simple --- una inundación de paquetes hacia un servicio --- de punta a punta:

\begin{enumerate}
\item El sensor conectado al espejo agrega el tráfico por destino y ventana, y publica muestras.
\item El monitor recibe las muestras y las publica para la detección.
\item La detección compara contra la línea base del destino: la desviación supera el umbral y publica una anomalía con su severidad propuesta.
\item El gestor de incidentes abre el incidente y lo publica.
\item El motor de políticas evalúa el incidente y decide: limitar la tasa del origen a 1000 paquetes por segundo, con vencimiento de 300 segundos.
\item La decisión viaja como orden al adaptador (por la API del controlador); su eco queda en el bus para auditoría y consola.
\item El adaptador instala el medidor y la regla de prioridad 1000, y confirma con barrera.
\item El evento \emph{mitigación aplicada} se publica con las reglas instaladas.
\item El monitor compara las tasas antes y después y publica \emph{mitigación verificada} (cumple o no el objetivo).
\item Al vencer los 300 segundos, la regla expira sola --- o se retira antes por cookie --- y el hecho se publica.
\item La auditoría registró cada paso, con su secuencia por incidente y sus huecos declarados si los hubo.
\end{enumerate}

\section{La elevación temporal}



El académico puede necesitar, por una práctica de laboratorio, un alcance que su perfil no tiene. El flujo completo:

\begin{enumerate}
\item El académico solicita la elevación por el portal: alcance y duración.
\item La solicitud se publica como evento; las políticas la registran como pendiente.
\item Un aprobador (distinto del solicitante --- condición C4) la aprueba; la aprobación queda vinculada a la solicitud y auditada.
\item El controlador instala la regla de elevación con prioridad 200 y vencimiento igual al concedido (1 hora por defecto).
\item Al vencer ---o al retirarse---, el retiro por cookie devuelve al equipo a su perfil académico, sin intervención.
\end{enumerate}

La elevación es una sesión temporal, no un rol nuevo: comparte cookie de sesión, prioridad y relojes; cambia el alcance y el tope.

\section{Secuencia completa con protocolo por salto}

El caso de uso completo ---conexión y login de un académico--- con los mensajes exactos entre componentes:

\begin{longtable}{p{1.2cm}p{4.8cm}p{4.4cm}p{3.6cm}}
\toprule
\textbf{Paso} & \textbf{Qué ocurre} & \textbf{Mensaje} & \textbf{Entre}\\
\midrule
1 & El equipo envía su primer paquete & Sin coincidencia $
ightarrow$ paquete al controlador & Switch $
ightarrow$ Controlador (OpenFlow)\\
2 & El controlador aprende (MAC, IP) y ubicación, e instala el esqueleto & Reglas 150/120/110/100/10 & Controlador $
ightarrow$ Switch\\
3 & La persona abre el portal y elige su reino & Formulario de credenciales & Persona $
ightarrow$ Portal (HTTPS)\\
4 & El portal reenvía las credenciales & Verificación & Portal $
ightarrow$ Identidad (REST)\\
5 & La identidad consulta al proveedor institucional & Petición de acceso RADIUS & Identidad $
ightarrow$ RADIUS (UDP)\\
6 & El proveedor responde & Acceso aceptado & RADIUS $
ightarrow$ Identidad\\
7 & La identidad crea la sesión y la publica & Evento \emph{sesión abierta} & Identidad $
ightarrow$ Bus (AMQP)\\
8 & La identidad ordena el perfil & Aplicar perfil (cookie de sesión) & Identidad $
ightarrow$ Controlador (REST)\\
9 & El controlador instala y confirma & Regla de prioridad 200 + barrera & Controlador $
ightarrow$ Switch\\
10 & El equipo navega por su alcance & Coincidencia en 200 $
ightarrow$ caminos & Switch $
ightarrow$ Switch (plano de datos)\\
11 & La auditoría registra cada paso & Eventos recibidos del catálogo & Bus $
ightarrow$ Auditoría\\
\bottomrule
\end{longtable}

El mismo patrón ---identificar, decidir, traducir, ejecutar, registrar--- se repite para la mitigación y para la elevación; lo único que cambia es quién dispara el primer paso.

\section{El bus: colas y enlaces}

Cada servicio consumidor tiene su cola propia, duradera, con sus enlaces al catálogo:

\begin{longtable}{p{3.2cm}p{6.4cm}p{4.4cm}}
\toprule
\textbf{Consumidor} & \textbf{Enlaces (tipos de evento)} & \textbf{Propósito}\\
\midrule
Auditoría & El catálogo completo & Registro total, sin excepción\\
Monitor & Dispositivo conectado \textperiodcentered{} mitigación aplicada & Ubicaciones y verificación\\
Detección & Muestra de métricas & La materia prima de la línea base\\
Incidentes & Anomalía \textperiodcentered{} verificación \textperiodcentered{} expiración \textperiodcentered{} dispositivo movido & Ciclo de vida del incidente\\
Políticas & Incidente abierto \textperiodcentered{} verificación \textperiodcentered{} expiración \textperiodcentered{} elevación \textperiodcentered{} dispositivo movido & Decisión y seguimiento\\
Identidad & Dispositivo movido & Invalidación de la sesión ligada\\
Registro & Dispositivo conectado & Historial por ubicación\\
Consola & Anomalía \textperiodcentered{} incidente abierto \textperiodcentered{} altas y bajas de dispositivos & Alertas en vivo\\
\bottomrule
\end{longtable}

Además, cada incidente abre su \textbf{cola propia}, declarada al abrirse y auto-borrada al cerrarse (con respaldo de 24 horas): es el mecanismo de entrega ordenada --- una sola cola, orden FIFO, sin reordenamiento en el consumidor. Los parámetros de entrega: confirmación de publicación activada; acuse manual tras procesar; cinco reintentos y luego cola de no entregados, visible y auditada; ventana de duplicados de 7 días por identificador; y un salto de secuencia se espera un instante y, si no llega, se declara el hueco --- jamás se procesa fuera de orden ni se oculta.

\section{Funciones internas de los servicios}

Las piezas internas que sostienen las secciones anteriores, servicio por servicio:

\begin{longtable}{p{3.4cm}p{4.4cm}p{6.2cm}}
\toprule
\textbf{Servicio} & \textbf{Función} & \textbf{Qué hace}\\
\midrule
Identidad & registrar\_operador & Crea la credencial (hash bcrypt) y el secreto TOTP; devuelve la URI del código QR, que se muestra una sola vez\\
Identidad & registrar\_invitado & Crea la cuenta efímera con su vencimiento\\
Identidad & verificar & Despacha por reino: RADIUS, repositorio propio con TOTP, o registro efímero\\
Identidad & emitir\_sesión & Crea la sesión con los relojes del perfil y publica el evento\\
Identidad & sesión\_válida & Aplica inactividad y tope al consultar; cierra y publica al vencer\\
Identidad & manejar\_evento & Un dispositivo movido invalida sus sesiones al instante\\
Identidad & purgar\_invitados & Bucle que borra las cuentas vencidas\\
Motor & decidir & Incidente abierto $
ightarrow$ acción y parámetros (tasa, TTL) según la política\\
Motor & enviar\_orden & La orden viaja al controlador; el eco queda en el bus\\
Adaptador & cookie\_hex & Clave legible $
ightarrow$ cookie de 64 bits con su tipo en los bits altos\\
Adaptador & aplicar\_mitigación & Medidor y regla 1000 según la acción; rechaza bloquear sin origen\\
Adaptador & aplicar\_perfil & Regla 200 con los dos relojes de la sesión\\
Adaptador & retirar & Elimina solo las reglas de la cookie propia\\
\bottomrule
\end{longtable}

El invariante del retiro, en pseudocódigo:

\begin{verbatim}
retirar(orden):
    cookie = registro_cookies[orden.contexto.cookie]
    si cookie no existe:
        responder 404 (cookie desconocida)
    para cada switch de la red:
        para cada regla instalada en el switch:
            si regla.cookie == cookie:
                eliminar regla
    responder retirada confirmada
\end{verbatim}

Ninguna otra función del sistema toca reglas ajenas: el retiro por cookie es la única vía, en todas las tablas.

\section{Análisis cuantitativo de R1 --- control de acceso según rol}



\begin{longtable}{p{6.4cm}p{2.2cm}p{5.4cm}}
\toprule
\textbf{Qué se mide} & \textbf{Cifra} & \textbf{Cómo se obtuvo}\\
\midrule
Roles definidos en la política & 4 & académico \textperiodcentered{} especialista de TI \textperiodcentered{} administrador de red \textperiodcentered{} superadministrador\\
Permisos y condiciones del modelo & 37 permisos \textperiodcentered{} 11 condiciones & Catálogo de la política, validado automáticamente\\
Asignaciones rol $\times$ permiso & 78 & Matriz generada y verificada contra su fuente\\
Reinos de identidad & 3 & comunidad (RADIUS) \textperiodcentered{} operadores (bcrypt + TOTP) \textperiodcentered{} invitados (TTL 24 h)\\
Entradas del perfil base por dispositivo & 6 & las cinco de la escalera más el aprendizaje\\
Sesión académica & 30 min de inactividad \textperiodcentered{} 12 h de tope & Valor inicial de diseño; se mide en campaña\\
Sesión de operador & 15 min de inactividad \textperiodcentered{} 8 h de tope & Más estricta: abre la gestión y la consola\\
Cookie de sesión y de incidente & bits altos 0x01 y 0x02 & Convención fija del diseño\\
Casos automatizados del módulo de identidad & 21 & Suite de pruebas en verde (verificación continua)\\
\bottomrule
\end{longtable}

Qué significa cada cifra, en una frase: cuatro roles bastan porque el privilegio fino sale de la combinación \emph{rol + recurso + condición}, no de más roles; setenta y ocho asignaciones cubren los treinta y siete permisos sin redundancias; los tres reinos garantizan que una cuenta de administrador nunca sea la misma que la de alumno; y los dos relojes por sesión aseguran que ningún privilegio quede abierto más de lo necesario.

\section{Análisis cuantitativo de R2 --- protección de recursos privilegiados}

\begin{longtable}{p{6.4cm}p{2.2cm}p{5.4cm}}
\toprule
\textbf{Qué se mide} & \textbf{Cifra} & \textbf{Cómo se obtuvo}\\
\midrule
Tipos de recurso catalogados & 50+ & El catálogo distingue servicios, recursos, cuentas, red y dispositivos\\
Clases mínimas de clasificación & 3 & generales \textperiodcentered{} privilegiados \textperiodcentered{} críticos\\
Denegación a la gestión desde el perfil base & descarte en prioridad 100 & Entrada permanente; ninguna sesión la abre\\
Cadena de cuentas privilegiadas & SA $
ightarrow$ Admin $
ightarrow$ TI & Delegación vigente, con responsable y vencimiento\\
Aprobación de alto impacto & obligatoria para cuarentena & Sin aprobación, la orden se rechaza en el controlador\\
Registro de dispositivos privilegiados & acotado & Solo operadores; los académicos no se registran\\
Accesos registrados & 100 \% & La auditoría escucha el catálogo completo\\
Retiro de reglas & siempre por cookie & Probado: una cookie ajena jamás se toca\\
\bottomrule
\end{longtable}

La traducción de la decisión a regla SDN quedó implementada y verificada: la mitigación se instala con prioridad 1000 y vencimiento igual a su TTL ---toda respuesta expira sola---, la sesión con prioridad 200, y el retiro por cookie está probado por caso automatizado. La clasificación del recurso decide qué condiciones lo protegen: un recurso crítico exige ventana de mantenimiento o emergencia declarada para su configuración.

\section{Análisis cuantitativo de RE --- eficiencia}

RE agrupa rendimiento, latencia, escalabilidad y continuidad. Las cifras consolidadas:

\begin{longtable}{p{5.6cm}p{2.6cm}p{5.8cm}}
\toprule
\textbf{Presupuesto} & \textbf{Cifra} & \textbf{Cómo se verifica}\\
\midrule
Entradas de flujo en la red (10 switches) & 420--520 nominales \textperiodcentered{} 665 peor caso & Presupuesto por componente; el llenado real se mide hasta agotar la tabla\\
Conmutación ante caída de enlace & grupo de respaldo o reinstalación & Se mide el tiempo en la campaña de pruebas\\
Confirmación de orden & barrera OpenFlow & La orden responde con la regla ya escrita\\
Reintento ante controlador caído & la misma orden, sin límite & Idempotencia por cookie\\
Bus: confirmación de publicación & activada & El emisor sabe que el mensaje fue recibido\\
Bus: reintentos y cola de no entregados & 5 reintentos \textperiodcentered{} cola por servicio & Lo no procesado queda visible y auditado\\
Ventana de duplicados & 7 días por identificador & Idempotencia del consumidor\\
\bottomrule
\end{longtable}

Recursos del despliegue ---el techo físico de RE en el prototipo---, 14 máquinas virtuales:

\begin{longtable}{p{5.6cm}p{2.4cm}p{2.4cm}p{3.2cm}}
\toprule
\textbf{Conjunto} & \textbf{Núcleos} & \textbf{RAM} & \textbf{Disco}\\
\midrule
4 switches virtuales & 4 & 4 GB & 32 GB\\
3 controladores (clúster) & 6 & 6 GB & 30 GB\\
4 servidores de plataforma & 8 & 18 GB & 59 GB\\
3 equipos de prueba & 3 & 3 GB & 24 GB\\
\textbf{Total} & \textbf{21} & \textbf{31 GB} & \textbf{145 GB}\\
\bottomrule
\end{longtable}

Peso en disco de los componentes de terceros tras su instalación: controlador SDN 559 MB, base de datos 158 MB, sensor de tráfico 135 MB, visualización 117 MB, bus 112 MB, métricas 102 MB, autenticación RADIUS 87 MB, directorio 87 MB; los once servicios propios rondan los 200 MB cada uno. El disco es la restricción dominante, y por eso las imágenes viajan empaquetadas, se cargan y se borran de cada máquina.

\section{Contraste con el avance previo y ajustes justificados}

El avance anterior declaró explícitamente no contener mediciones; esta consolidación cierra esa deuda y corrige lo que la evidencia corrigió:

\begin{longtable}{p{5.0cm}p{4.6cm}p{4.4cm}}
\toprule
\textbf{Supuesto del avance previo} & \textbf{Estado consolidado} & \textbf{Justificación}\\
\midrule
Diez switches dedicados & Cuatro switches virtuales (mínimo tres) & Los recursos entregados son 14 máquinas; se conserva núcleo dual y acceso con doble enlace\\
Memoria TCAM sin verificar & Sin TCAM: conmutación por software & La sonda lo confirmó; el límite se mide por llenado\\
Espejo de tráfico en el acceso & Espejo en el núcleo, entrega por confirmar & El núcleo ve el tráfico que interesa a la detección; hay respaldo fijado\\
Red de gestión numerada con .0 & VLAN 100 = 10.0.0.0/24, canon 10/20/30/100 & Reordenamiento cerrado en equipo\\
Múltiples tablas en duda & 254 tablas confirmadas & Sonda directa\\
Clúster de dos o tres nodos & Tres nodos & El consenso exige mayoría\\
Servicios «por diseñar» & Once servicios construidos y probados & 186 casos automatizados en verde\\
\bottomrule
\end{longtable}

\section{Cumplimiento de R1, R2 y RE}

\begin{longtable}{p{2.4cm}p{6.0cm}p{5.6cm}}
\toprule
\textbf{Requisito} & \textbf{Dónde se cumple} & \textbf{Evidencia}\\
\midrule
R1.1--R1.3 & Identificación por MAC y contexto; autenticación por reino; perfil por defecto y elevado & Portal + identidad; sesiones con perfil y casos automatizados\\
R1.4--R1.6 & Política contextual y traducción a reglas; mínimo privilegio; denegación & 37 permisos, 78 asignaciones, descarte a gestión, denegación por defecto\\
R1.7--R1.8 & Gestión de roles, perfiles y dispositivos; rol administrativo restringido & Cadena SA$
ightarrow$Admin$
ightarrow$TI; segundo factor obligatorio\\
R1.9--R1.10 & Trazabilidad y crecimiento acotado & Auditoría completa; registro solo privilegiado\\
R2.1--R2.9 & Inventario, clasificación, políticas, segmentación, actualización y auditoría & Catálogo de 50+ tipos, tres clases, descarte a gestión, aprobación de alto impacto\\
R2.10 & Traducción de decisión a regla SDN & Adaptador: prioridades 1000/200, cookies, retiro por cookie\\
RE & Rendimiento, latencia, escalabilidad y continuidad & Presupuestos y recursos de §24; medición fina en campaña\\
\bottomrule
\end{longtable}

\section{Declaración de uso de inteligencia artificial}
Conforme a los lineamientos del curso, el Grupo 3 declara que usó inteligencia artificial generativa (Claude, de Anthropic, a través de Claude Code) para consolidar el diseño con la implementación, extraer las cifras cuantitativas de las evidencias de laboratorio y redactar este documento. Las decisiones de diseño, la interpretación de los resultados y la responsabilidad por el contenido son del equipo, que revisó y valida cada cifra aquí presentada.

\end{document}
